\documentclass[11pt]{article}
\usepackage[T1]{fontenc}
\usepackage[utf8]{inputenc}
\usepackage[margin=1in]{geometry}
\usepackage{booktabs}
\usepackage{xcolor}
\usepackage{listings}
\usepackage{enumitem}
\usepackage{parskip}
\usepackage[colorlinks=true,allcolors=blue!55!black]{hyperref}

\lstset{
  basicstyle=\ttfamily\small,
  backgroundcolor=\color{black!4},
  frame=single,
  breaklines=true,
  breakatwhitespace=true,
  columns=fullflexible,
}

\title{\textbf{GAMECA} v0.4.39\\[4pt]
       \large Installation, Quick-Start \& Output Reference}
\author{Anmol Dash}
\date{\today}

\begin{document}
\maketitle
\tableofcontents
\newpage

% ===========================================================================
\section{What is GAMECA?}

\textbf{GAMECA} (Gene Alignment, Motif, Expression and Clustering Analysis)
is a native macOS desktop application for the comprehensive downstream
characterisation of transposable-element (TE) families. It ships as a
\texttt{.dmg} installer containing a Tauri/React frontend and a bundled
Python analysis backend. Starting from nothing more than a RepeatMasker
family name and a genome assembly code, GAMECA runs the following stages:

\begin{enumerate}[itemsep=2pt]
  \item \textbf{Data prep} --- fetch all annotated genomic copies from
        UCSC RepeatMasker or Dfam, then extract sequences (from a local
        genome FASTA or via the UCSC API).
  \item \textbf{Clustering} --- group copies by $k$-mer profile
        (SVD \textrightarrow\ UMAP/t-SNE \textrightarrow\ HDBSCAN).
  \item \textbf{Alignment} --- per-cluster MAFFT alignment, CIAlign
        cleaning, and majority-vote consensus sequences.
  \item \textbf{Motif analysis} --- scan every copy for JASPAR
        transcription-factor binding motifs.
  \item \textbf{Enrichment} --- Fisher's exact test for motifs
        over-represented in each cluster.
  \item \textbf{TFBS analysis} --- count and visualise transcription-factor
        binding sites per cluster and per strand.
  \item \textbf{GO annotation} --- map enriched motifs to their
        transcription factors and Gene Ontology terms.
  \item \textbf{Expression} --- per-cluster expression boxplots and
        differential-expression tests (when an expression matrix is
        provided).
  \item \textbf{Primer design} --- family- and cluster-specific PCR
        primers with genome-wide specificity checking.
  \item \textbf{Stage~11 standout analyses} --- phylogenetics, CRISPR
        gRNA off-target scoring, 3$'$ transduction detection, antisense
        promoter scanning, CTCF/TAD overlap, epigenetic overlay, ortholog
        insertion presence, multi-assembly liftover, ColabFold consensus
        ORF structure prediction, repeat landscape, LTR structural
        classification, subfamily resolution, and differential expression
        between clusters.
\end{enumerate}

Analyses run either \textbf{locally} on your Mac or are submitted as
independent \textbf{batch jobs} on an LSF or Slurm HPC cluster over SSH.
No terminal knowledge is required for typical use.

% ===========================================================================
\section{System requirements}

\begin{tabular}{@{}lp{0.62\linewidth}@{}}
\toprule
Component & Requirement \\
\midrule
macOS & 12 Monterey or later (Apple Silicon \emph{or} Intel) \\
Disk  & $\sim$3~GB (app + Python environment + first ColabFold model weights) \\
RAM   & 8~GB minimum; 16~GB or more recommended for large families \\
Internet & Required for UCSC/Dfam sequence fetching, first-run package
           install, JASPAR/GO lookups, and HPC SSH \\
HPC cluster & Optional; LSF (\texttt{bsub}) or Slurm (\texttt{sbatch})
              reachable by SSH from your Mac \\
\bottomrule
\end{tabular}

% ===========================================================================
\section{Download}

Installers are published on the GitHub Releases page:

\begin{center}
  \url{https://github.com/anmol-dash/te-scraper-and-analysis/releases}
\end{center}

Download the \texttt{.dmg} that matches your Mac's processor:

\begin{tabular}{@{}ll@{}}
\toprule
Architecture & File to download \\
\midrule
Apple Silicon (M1/M2/M3/M4) & \texttt{GAMECA\_0.4.39\_aarch64.dmg} \\
Intel & \texttt{GAMECA\_0.4.39\_x64.dmg} \\
\bottomrule
\end{tabular}

\smallskip
If unsure which processor your Mac has, click the Apple menu
\textrightarrow\ \textbf{About This Mac} and read the ``Chip'' or
``Processor'' line.

% ===========================================================================
\section{Installation}

\begin{enumerate}
  \item Double-click the downloaded \texttt{.dmg} to mount it.
  \item Drag \textbf{GAMECA.app} into your \textbf{Applications} folder.
  \item Eject the disk image (\texttt{Cmd-E}).
  \item \textbf{First launch only --- Gatekeeper:} because GAMECA is
        distributed outside the Mac App Store, macOS blocks the first
        open. \textbf{Right-click} (or Control-click) \textbf{GAMECA.app}
        \textrightarrow\ \textbf{Open} \textrightarrow\ \textbf{Open}.
        This prompt appears once only.
\end{enumerate}

% ===========================================================================
\section{First launch and automatic setup}

On first launch GAMECA performs a one-time setup ($\sim$2--3 minutes):

\begin{enumerate}
  \item The bundled Python sidecar starts and sends a \texttt{ping}
        handshake to the frontend; a spinner shows until it completes.
  \item The app checks for a Python virtual environment at
        \texttt{\$HOME/gameca\_venv}. If absent, it creates one and
        installs all dependencies from the bundled \texttt{requirements.txt}
        (NumPy, pandas, scikit-learn, UMAP, HDBSCAN, BioPython, MAFFT
        bindings, CIAlign, MOODS, \texttt{colabfold[alphafold]}, and more).
  \item A progress bar and log stream in the \textbf{Log} panel. Do not
        close the app during this step.
  \item When the status indicator turns green the app is ready.
\end{enumerate}

\noindent\textbf{Note:} ColabFold model weights ($\sim$4~GB) are
\emph{not} downloaded here; they are fetched the first time a
structure-prediction job runs and cached in
\texttt{\$HOME/.cache/colabfold}.

% ===========================================================================
\section{The interface}

\begin{description}[leftmargin=0pt]
  \item[Sidebar (left).] Every pipeline step, grouped into \emph{Data
        prep}, \emph{Core analysis}, \emph{Enrichment}, \emph{Primer
        design}, \emph{Full pipeline}, and \emph{Results}. Click a step
        to open its configuration panel.
  \item[Configuration panel (centre).] Parameters for the selected step:
        family name, assembly, output directory, queue, and step-specific
        options. Fill in and click \textbf{Run on Cluster} or
        \textbf{Run locally}.
  \item[Log panel (bottom).] Streams live output line-by-line: stage
        checkpoints, progress bars, and errors.
\end{description}

% ===========================================================================
\section{Connecting to an HPC cluster}

\begin{enumerate}
  \item Click \textbf{Settings} (gear icon) \textrightarrow\
        \textbf{HPC Connection}.
  \item Enter \textbf{Hostname}, \textbf{Username}, \textbf{SSH key path}
        (blank to use the SSH agent), \textbf{Scheduler}
        (\texttt{LSF}/\texttt{Slurm}/\emph{Auto}),
        \textbf{Queue / partition}, and a writable \textbf{Work
        directory} (e.g.\ \texttt{/home/amodz/anmol/gameca}).
  \item Click \textbf{Test connection}. A green tick confirms the
        scheduler was detected and the work directory is writable.
\end{enumerate}

\noindent\textbf{Duo / two-factor clusters:} if your cluster requires MFA
on every login, add a \texttt{ControlMaster} entry to
\texttt{\textasciitilde/.ssh/config} so the app reuses one authenticated
session:

\begin{lstlisting}
Host hpclogin1.university.edu
    ControlMaster auto
    ControlPath  ~/.ssh/cm-%r@%h:%p
    ControlPersist 8h
\end{lstlisting}

% ===========================================================================
\section{Running an analysis}

\subsection{Check the locus count first}

\textbf{Data prep} \textrightarrow\ \textbf{RMSK locus count} (or
\textbf{Dfam locus count}) verifies the family name and reports how many
copies exist before you commit to a full job.

\subsection{Single family --- on the cluster}

\begin{enumerate}
  \item \textbf{Full pipeline} \textrightarrow\ \textbf{Batch job}.
  \item Set \textbf{Family} (e.g.\ \texttt{THE1A}), \textbf{Assembly}
        (\texttt{hg38}/\texttt{hg19}/\texttt{mm10}/\texttt{mm39}),
        \textbf{Output directory}, and optionally \textbf{Notify email}.
  \item Click \textbf{Run on Cluster}. GAMECA uploads the scripts, writes
        a \texttt{bsub}/\texttt{sbatch} job, and submits it. The job ID
        prints in the Log panel.
  \item Monitor in \textbf{Results} \textrightarrow\ \textbf{File browser}
        or with \texttt{bjobs -w} on the cluster.
  \item When done, \textbf{Results} \textrightarrow\ \textbf{Retrieve}
        \texttt{rsync}s the output folder back to your Mac.
\end{enumerate}

\subsection{Single family --- locally}

\textbf{Full pipeline} \textrightarrow\ \textbf{Batch job}, fill in
\textbf{Family} and \textbf{Assembly}, then \textbf{Run locally}. Output
lands in \texttt{\$HOME/gameca\_results/<family>/}. Suitable for families
up to $\sim$5{,}000 copies on 16~GB RAM; larger families and ColabFold
prediction should run on HPC.

\subsection{Many families --- batch submission}

Submit many families as independent jobs with the bundled
\texttt{submit\_diego\_families.py}. Each family produces \emph{two}
separate LSF jobs --- a main job (core pipeline + Stage~11 except
ColabFold) and a dependent fold job (ColabFold).

\begin{lstlisting}
source ~/gameca_venv/bin/activate

python submit_diego_families.py \
    --outdir /home/$USER/gameca_results \
    --queue  normal \
    --notify-email you@institution.edu

python submit_diego_families.py --dry-run   # preview without submitting
\end{lstlisting}

Both jobs per family appear immediately in \texttt{bjobs}:
\texttt{gam\_<FAMILY>\_main} starts when a slot frees;
\texttt{gam\_<FAMILY>\_fold} stays \textsc{pend} until its main job
finishes, then starts automatically. Two emails are sent per family
(one when the main job ends, one when the fold job ends).

% ===========================================================================
\section{Output reference}

A run writes one folder per family. The tree below is the complete
layout; the subsections that follow show \emph{what is inside each file}
so you can read the results without ever opening the app.

\begin{lstlisting}
<outdir>/<family>/                  e.g. gameca_results/the1a/
  01_data/
    the1a_clustered.csv             master table: every copy + cluster
  02_statistics/
    cluster_summary.csv             one row per cluster (size, GC, length)
    per_cluster/cluster_<c>_statistics.txt
  03_clustering/
    clustering_visualization.html   interactive UMAP scatter
    clustering_visualization_expr.html   (if expression provided)
  04_alignments/
    alignment_stats.txt
  05_consensus/
    the1a_consensus.fa              one consensus per cluster + global
  05_tfbs/
    cluster_tf_binding_counts.csv   TF binding sites per cluster
    strand_tf_binding_counts.csv
    cluster_tf_binding_heatmap.html
    tf_binding_top_overlaps.html
  06_primers/
    cluster_top5_primers.csv        recommended primers per cluster
    selected_primers_summary.csv
  07_visualizations/
    index.html                      dashboard linking every figure
  cluster_alignments/
    cluster_<c>_seqs.fa / _aligned.fa / _consensus.fa
    all_cluster_consensuses.fa
  cialign_plots/
    index.html                      open in a browser to view alignments
  motif_analysis/
    all_overlaps.tsv                every copy x motif overlap
    overall_motif_counts.csv        motif frequency table
    overall_top_motifs.png
  enrichment_results/
    cluster_<c>_enrichment.csv      Fisher test per motif per cluster
    cluster_<c>_enrichment_annotated.csv
    enrichment_heatmap.png
  go_annotations/
    gene_functions.csv              TF -> gene -> GO lookups
    go_bp_top_terms.csv
    go_bp_top_terms.png
  reports/                          Stage 11 results (see 9.7)
    standout_analysis.log
    fig_*.pdf / *_per_copy.csv / *_report.txt / *_measured_values.tex
    fold/                           ColabFold 3D structures
  bsub_main.out / bsub_fold.out     live cluster job logs
\end{lstlisting}

% ---------------------------------------------------------------------------
\subsection{Clustering: the master table}

\textbf{\texttt{01\_data/<family>\_clustered.csv}} is the central output:
one row per genomic copy. \texttt{Cluster} is the group it was sorted
into; \texttt{chr/start/stop/strand} give its location; \texttt{Seq} is
its DNA; \texttt{umap\_x/umap\_y} are the 2-D coordinates plotted in
\texttt{03\_clustering/clustering\_visualization.html}.

\begin{lstlisting}
Seq,Cluster,chr,start,stop,strand,te_name,umap_x,umap_y
tgttattagacgcg...,3,chr1,3031358,3031710,-,IAPLTR1a_Mm,4.21,-1.08
aatattcagttgtc...,1,chr2,9881204,9881559,+,IAPLTR1a_Mm,-2.66,3.94
\end{lstlisting}

\textbf{\texttt{02\_statistics/cluster\_summary.csv}} summarises each
cluster (number of copies, mean length, GC content). The interactive
UMAP scatter in \texttt{03\_clustering/} colours every copy by cluster ---
open it in a browser and hover for per-copy detail.

% ---------------------------------------------------------------------------
\subsection{Alignment and consensus}

\textbf{\texttt{cluster\_alignments/cluster\_<c>\_consensus.fa}} is the
``average'' sequence of each cluster in FASTA format; the \texttt{>} line
is the name and the lines below are the sequence.
\textbf{\texttt{all\_cluster\_consensuses.fa}} bundles them all and is
the input for downstream structure/divergence modules.

\begin{lstlisting}
>IAP_cluster3_consensus
GGGGCTATTTTTGGTTGGTGTGGACGTTTTTTGTCTGGAAAATTAAGAGG
ATGCCAAAAGGTCCGATTTGAAAAAACAGTCTTTACATTGGAAATATCTA
\end{lstlisting}

\textbf{\texttt{cialign\_plots/index.html}} is an HTML gallery of the
CIAlign before/after/markup images --- the quickest way to judge how
clean each cluster's alignment is.

% ---------------------------------------------------------------------------
\subsection{Motif analysis (\texttt{motif\_analysis/})}

This stage scans every copy for JASPAR transcription-factor binding
motifs. \textbf{\texttt{all\_overlaps.tsv}} is the raw result: one line
per (copy, motif) overlap, giving the copy coordinates, the motif name,
its position, score, and strand --- this is the table every downstream
enrichment/TFBS step is built from.

\textbf{\texttt{overall\_motif\_counts.csv}} collapses that to a simple
frequency table (how many times each motif was found across all copies),
and \textbf{\texttt{overall\_top\_motifs.png}} plots the top~20.

\begin{lstlisting}
Motif_name,count
CTCF,1840
SP1,1422
KLF4,1190
\end{lstlisting}

% ---------------------------------------------------------------------------
\subsection{Enrichment (\texttt{enrichment\_results/})}

Finding a motif is not the same as it being \emph{characteristic} of a
cluster. For each cluster, GAMECA runs a Fisher's exact test asking
``is this motif over-represented in this cluster versus the rest of the
family?'' \textbf{\texttt{cluster\_<c>\_enrichment.csv}} holds one row
per motif, sorted by \texttt{P\_Value}:

\begin{lstlisting}
Motif,In_Cluster,Total_With_Motif,Cluster_Size,Total_Loci,Odds_Ratio,P_Value
CTCF,118,402,140,1485,6.91,3.2e-41
SP1,77,388,140,1485,3.10,1.1e-12
\end{lstlisting}

\begin{itemize}[itemsep=1pt]
  \item \texttt{In\_Cluster} --- copies in this cluster with the motif.
  \item \texttt{Total\_With\_Motif} --- copies family-wide with the motif.
  \item \texttt{Odds\_Ratio} --- how much more frequent the motif is in
        this cluster ($>1$ = enriched).
  \item \texttt{P\_Value} --- Fisher's exact significance; rows with
        \texttt{P\_Value} below the threshold are the cluster's
        signature TFs.
\end{itemize}

\textbf{\texttt{cluster\_<c>\_enrichment\_annotated.csv}} adds the
transcription-factor name, a one-line summary, and GO terms (see
\S9.6) to every significant motif. \textbf{\texttt{enrichment\_heatmap.png}}
shows clusters $\times$ enriched motifs at a glance.

% ---------------------------------------------------------------------------
\subsection{TFBS analysis (\texttt{05\_tfbs/})}

Where enrichment asks ``which motifs are statistically special'', the
TFBS stage simply \emph{counts} binding sites so you can compare raw
binding load across clusters.
\textbf{\texttt{cluster\_tf\_binding\_counts.csv}} has one row per
(cluster, motif):

\begin{lstlisting}
Cluster,Motif_name,Binding_Sites
0,CTCF,118
0,SP1,77
1,REST,54
\end{lstlisting}

\textbf{\texttt{strand\_tf\_binding\_counts.csv}} breaks the same counts
down by DNA strand. The two \texttt{.html} files are interactive Plotly
views: \texttt{cluster\_tf\_binding\_heatmap.html} (clusters $\times$ TFs,
log-scaled) and \texttt{tf\_binding\_top\_overlaps.html} (the 30 highest
single overlaps). Open either in a browser.

% ---------------------------------------------------------------------------
\subsection{GO annotation (\texttt{go\_annotations/})}

This stage answers ``what do the enriched transcription factors actually
\emph{do}?'' Each significant motif is mapped to its gene via
\texttt{mygene.info}, and \textbf{\texttt{gene\_functions.csv}} records
the function and Gene Ontology terms:

\begin{lstlisting}
Gene,Symbol,Name,Summary,GO_BP,GO_MF,GO_CC
CTCF,CTCF,CCCTC-binding factor,"Chromatin insulator...",chromatin organization;...,DNA binding;...,nucleus;...
\end{lstlisting}

\textbf{\texttt{go\_bp\_top\_terms.csv}} ranks the most common Biological
Process terms across all enriched TFs, and \texttt{go\_bp\_top\_terms.png}
(plus per-cluster \texttt{cluster\_<c>\_go\_bp.png}) plots them.

\begin{lstlisting}
GO_Term,Count
regulation of transcription,42
chromatin organization,28
\end{lstlisting}

\noindent If no motifs pass the significance threshold the stage writes a
\texttt{GO\_SKIPPED.txt} note instead of empty files --- expected for
small or highly divergent families.

% ---------------------------------------------------------------------------
\subsection{Primers (\texttt{06\_primers/})}

\textbf{\texttt{cluster\_top5\_primers.csv}} lists the recommended PCR
primers per cluster. \texttt{coverage} is how many copies the primer
matches; \texttt{total\_expr} sums their expression; \texttt{rows\_covered}
names the copies. \texttt{selected\_primers\_summary.csv} is the global
shortlist with the selection \texttt{strategy} used.

\begin{lstlisting}
cluster,primer,coverage,total_expr,rows_covered
0,CCTAGCTATGCTGCCTTG,6,1119.19,"8,9,10,11,12,13"
\end{lstlisting}

% ---------------------------------------------------------------------------
\subsection{Stage 11 standout analyses (\texttt{reports/})}

Every Stage~11 analysis writes the \emph{same four kinds of file}, so once
you can read one you can read them all. Using phylogenetics
(\texttt{phylo}) as the example:

\begin{tabular}{@{}lp{0.58\linewidth}@{}}
\toprule
File & What it is \\
\midrule
\texttt{fig\_phylo\_tree.pdf}
  & The figure --- open and look at it. Each analysis makes one or more
    \texttt{fig\_*} PDFs/PNGs. \\[4pt]
\texttt{phylo\_per\_copy.csv}
  & One row per copy, adding this analysis's columns (here: age,
    divergence, ORF status). Join back to the master table by
    coordinates or \texttt{TE\_ID}. \\[4pt]
\texttt{phylo\_report.txt}
  & A plain-English summary of the headline numbers. \\[4pt]
\texttt{phylo\_measured\_values.tex}
  & The same numbers as \LaTeX\ macros for the PDF report; ignore unless
    compiling the paper. \\
\bottomrule
\end{tabular}

\medskip
The \texttt{*\_report.txt} files need no software and are the fastest way
to read a result:

\begin{lstlisting}
GAMECA Phylogenetics / Divergence-Age --- Measured Values
  Family:                IAP
  Copies analysed:       12577
  Median est. age:       423.18 Myr
  Intact (ORF>=thr):     3374
  -- Top master/source element --
  Name:                  chr8:4886476
\end{lstlisting}

\noindent Some analyses also write a specialised table, e.g.\
\texttt{grna\_offtarget\_guides.csv} (CRISPR guides with coverage,
off-target count, specificity), \texttt{ltr\_struct\_summary.csv}
(full-length vs solo LTR counts), \texttt{subfamily\_tree.nwk} (a Newick
tree for FigTree/iTOL), and \texttt{divergence\_per\_locus.csv} (per-copy
divergence for the repeat-landscape plot).

\paragraph{Structure predictions (\texttt{reports/fold/}).}
If ColabFold ran, this folder holds the predicted 3-D structure of each
cluster's consensus ORF: \texttt{*.pdb} (open in PyMOL, ChimeraX, or
\url{https://molstar.org/viewer}), \texttt{*\_plddt.png} (per-residue
confidence; $>70$ reliable, $>90$ high), and \texttt{*\_pae.png}
(predicted aligned error; blue blocks = confident domains).

% ---------------------------------------------------------------------------
\subsection{Following a job and the dashboard}

\texttt{07\_visualizations/index.html} is a dashboard linking every
figure produced. \texttt{bsub\_main.out} / \texttt{bsub\_fold.out} are the
live cluster logs, and \texttt{reports/standout\_analysis.log} records
each Stage~11 step with timings --- check it first if a figure is missing:

\begin{lstlisting}
[phylo]      ... ok (612s)
[grna]       ... ok (94s)
[fold]       ... FAILED rc=1 (12s)   <- look here if a result is absent
\end{lstlisting}

\medskip
\noindent To compile the auto-generated PDF report stitching every figure
and number together:

\begin{lstlisting}
tectonic gameca_report.tex   # or: pdflatex gameca_report.tex
\end{lstlisting}

% ===========================================================================
\section{Updates}

\begin{description}[leftmargin=0pt]
  \item[Python scripts] are refreshed silently from the latest commit on
        \texttt{main} every launch. No action required.
  \item[App binary] When a new GitHub Release has a higher version, a
        banner appears on next launch. Click \textbf{Install update}; the
        new \texttt{.dmg} is downloaded, signature-verified, and
        installed. Relaunch when prompted.
\end{description}

To update manually, download the latest \texttt{.dmg} from Releases and
repeat \S4.

% ===========================================================================
\section{Troubleshooting}

\begin{description}[leftmargin=0pt,itemsep=8pt]

  \item[\textit{``GAMECA is damaged and cannot be opened.''}]
        macOS quarantine. Right-click \textbf{GAMECA.app} \textrightarrow\
        \textbf{Open} \textrightarrow\ \textbf{Open} (one time). If it
        persists, \textbf{System Settings} \textrightarrow\ \textbf{Privacy
        \& Security} \textrightarrow\ \textbf{Open Anyway}.

  \item[\textit{Blank screen or endless spinner.}]
        The Python sidecar failed to start. Confirm you downloaded the
        correct architecture. Open \textbf{Console.app}, filter for
        ``GAMECA'', and look for \texttt{ImportError} / \texttt{sidecar
        not found}.

  \item[\textit{Dependency install fails on first launch.}]
        Run it manually, then relaunch:
\begin{lstlisting}
python3 -m venv --system-site-packages ~/gameca_venv
source ~/gameca_venv/bin/activate
pip install -r \
  /Applications/GAMECA.app/Contents/Resources/requirements.txt
\end{lstlisting}

  \item[\textit{HPC connection fails.}]
        Verify SSH works from Terminal
        (\texttt{ssh -i \textasciitilde/.ssh/id\_ed25519 user@host}). For
        Duo clusters add a \texttt{ControlMaster} entry (see \S8).

  \item[\textit{Jobs stuck in \textsc{pend}.}]
        Normal when the queue is full or a fold job is waiting on its main
        job. Check slots with \texttt{bqueues normal}; inspect a specific
        job with \texttt{bjobs -l <JOBID>}.

  \item[\textit{GO stage produced only \texttt{GO\_SKIPPED.txt}.}]
        No motifs passed the enrichment p-value threshold --- expected for
        small or highly divergent families. Lower \texttt{--p-threshold}
        or inspect \texttt{enrichment\_results/} directly.

  \item[\textit{Motif/TFBS folders are empty.}]
        JASPAR motif data could not be fetched (network/proxy). Check
        \texttt{echo \$https\_proxy}; GAMECA reads standard proxy env
        vars. A \texttt{MOTIF\_ANALYSIS\_FAILED.txt} marker explains why.

  \item[\textit{ColabFold hangs or crashes on HPC.}]
        Old host GCC. Use a Singularity image:
\begin{lstlisting}
python run_stage11_all.py --only fold \
    --singularity-image /path/to/colabfold.sif \
    --no-gpu   # omit if GPUs are available
\end{lstlisting}

  \item[\textit{Sequences are all ``N'' or empty.}]
        The UCSC API was unreachable from the compute node. Check the
        proxy as above, or supply a local genome FASTA via
        \texttt{--genome}.

\end{description}

\end{document}
